\section{Introduction}

La conception permet de passer des besoins fonctionnels à une organisation technique claire. Dans ce
chapitre, nous présentons l'architecture générale de PharmacieConnect, la structure de chaque composant,
le modèle de données, les diagrammes UML prévus, les mécanismes de sécurité et les choix techniques
retenus.

\section{Architecture globale de la solution}

PharmacieConnect repose sur une architecture client-serveur. L'application mobile et le tableau de bord
administrateur ne communiquent pas directement avec la base de données. Toutes les opérations passent
par l'API backend. Ce choix permet de centraliser la validation, la sécurité et les règles métier.

\placeholderfigure
{Architecture globale de PharmacieConnect}
{Backend FastAPI, dashboard administrateur React/Vite, application mobile Expo/React Native et base de données SQLAlchemy avec SQLite ou PostgreSQL.}
{fig:architecture-globale}

L'application mobile consomme surtout les endpoints publics et les services liés au compte utilisateur. Le
dashboard administrateur utilise les endpoints protégés pour gérer les données. Le backend sert donc de
couche commune entre les deux interfaces et la base de données.

\section{Architecture du backend}

Le backend est développé avec FastAPI. Sa structure sépare les routes HTTP, les services métier, les
modèles de données et les schémas de validation. Cette séparation rend le code plus lisible et facilite
les tests.

\begin{table}[H]
    \centering
    \caption{Organisation du backend}
    \begin{tabular}{L{4cm}L{9cm}}
        \toprule
        \textbf{Élément} & \textbf{Rôle} \\
        \midrule
        \texttt{main.py} & Initialisation de l'application FastAPI, middlewares, migrations au démarrage et endpoints publics. \\
        \texttt{routers/} & Routes d'authentification, d'administration et d'analytics. \\
        \texttt{services/} & Logique métier liée à l'authentification, aux pharmacies, aux gardes, aux médicaments, à l'audit et au cache. \\
        \texttt{models.py} & Modèles SQLAlchemy représentant les tables principales. \\
        \texttt{schemas.py} & Schémas Pydantic utilisés pour valider les entrées et structurer les sorties. \\
        \texttt{migrations/} & Scripts SQL permettant de faire évoluer le schéma de la base. \\
        \texttt{events/} & Bus d'événements interne et écouteurs associés. \\
        \bottomrule
    \end{tabular}
\end{table}

Le traitement d'une requête suit généralement le même parcours. Le client appelle une route FastAPI, la
route vérifie les paramètres et l'utilisateur courant, puis le service métier applique les règles nécessaires
avant l'accès à la base via SQLAlchemy.

\section{Architecture de l'interface d'administration}

L'interface d'administration est développée avec React et Vite. Elle regroupe des pages, des composants
réutilisables, des contextes applicatifs et un client API centralisé.

\begin{table}[H]
    \centering
    \caption{Organisation de l'administration web}
    \begin{tabular}{L{4cm}L{9cm}}
        \toprule
        \textbf{Élément} & \textbf{Rôle} \\
        \midrule
        \texttt{src/pages/} & Pages principales du tableau de bord : pharmacies, gardes, médicaments, carte, audit, paramètres et gestion. \\
        \texttt{src/components/} & Composants d'interface, tableaux, formulaires, modales, cartes statistiques et éléments de layout. \\
        \texttt{src/context/} & Gestion de l'authentification, de la langue et de certains états globaux. \\
        \texttt{src/lib/api.js} & Client Axios responsable des appels backend et de l'injection du jeton JWT. \\
        \texttt{src/styles/} & Styles et utilitaires liés à la présentation de l'interface. \\
        \bottomrule
    \end{tabular}
\end{table}

Cette organisation évite de mélanger l'affichage avec la logique d'accès aux données. Le client API ajoute
le jeton d'accès aux requêtes protégées et gère le renouvellement de session lorsque le jeton expire.

\section{Architecture de l'application mobile}

L'application mobile est développée avec Expo et React Native. Elle est organisée autour d'écrans
orientés vers l'usage citoyen : recherche, carte, calendrier de garde, médicaments, authentification et
paramètres.

\begin{table}[H]
    \centering
    \caption{Organisation de l'application mobile}
    \begin{tabular}{L{4cm}L{9cm}}
        \toprule
        \textbf{Élément} & \textbf{Rôle} \\
        \midrule
        \texttt{screens/} & Écrans de l'application : accueil, carte, calendrier, médicaments, connexion, inscription, vérification email et paramètres. \\
        \texttt{components/} & Composants réutilisables comme les cartes de pharmacies, modales, sélecteurs et éléments RTL. \\
        \texttt{config/api.js} & Configuration de l'URL backend et des appels API. \\
        \texttt{utils/} & Fonctions utilitaires pour les données, les couleurs, le thème, le logging et la gestion RTL. \\
        \texttt{locales/} & Fichiers de traduction français, anglais et arabe. \\
        \bottomrule
    \end{tabular}
\end{table}

Des contextes React Native sont utilisés pour partager certaines informations entre les écrans, comme
l'authentification, la langue, le thème, les favoris, les filtres et les notifications. Cela évite de répéter
la même logique dans plusieurs parties de l'application.

\section{Modèle de données}

Le modèle de données est défini avec SQLAlchemy. Les principales entités utilisées par le backend sont
présentées dans le tableau suivant.

\begin{longtable}{L{4cm}L{10cm}}
    \caption{Entités principales du modèle de données}\\
    \toprule
    \textbf{Entité} & \textbf{Description} \\
    \midrule
    \endfirsthead
    \toprule
    \textbf{Entité} & \textbf{Description} \\
    \midrule
    \endhead
    \texttt{Utilisateur} & Compte utilisateur de l'application mobile avec email, mot de passe, statut actif et vérification email. \\
    \texttt{Administrateur} & Compte de gestion avec rôle, statut, informations de profil et périmètre régional éventuel. \\
    \texttt{RefreshToken} & Jeton de rafraîchissement permettant de prolonger une session selon les règles de sécurité. \\
    \texttt{LoginAttempt} & Tentative de connexion utilisée pour suivre les succès et les échecs d'authentification. \\
    \texttt{Pharmacie} & Données d'une pharmacie : nom, adresse, téléphone, gouvernorat, latitude, longitude et référence OSM éventuelle. \\
    \texttt{GardeSchedule} & Planning de garde avec date, pharmacie, horaire, ville, gouvernorat, type de garde et notes. \\
    \texttt{Medicine} & Catalogue de médicaments avec code PCT, nom commercial, prix public, tarif de référence, remboursement, DCI et AP. \\
    \texttt{AuditLog} & Trace des actions importantes avec acteur, cible, statut, adresse IP et détails. \\
    \texttt{SearchEvent} & Événements de recherche utilisés pour alimenter les statistiques et les analytics. \\
    \bottomrule
\end{longtable}

\section{Diagrammes UML}

Les diagrammes UML ci-dessous sont prévus pour compléter la conception. Les emplacements sont gardés
dans le rapport afin de pouvoir insérer les exports finaux dès qu'ils sont disponibles.

\placeholderfigure
{Diagramme de classes}
{Image UML à insérer : classes Utilisateur, Administrateur, Pharmacie, GardeSchedule, Medicine, AuditLog, RefreshToken, LoginAttempt et SearchEvent.}
{fig:class-diagram-placeholder}

\placeholderfigure
{Diagramme de séquence : authentification JWT}
{Image UML à insérer : saisie des identifiants, appel API, vérification du mot de passe, génération du jeton d'accès et du jeton de rafraîchissement.}
{fig:sequence-auth-placeholder}

\placeholderfigure
{Diagramme de séquence : recherche d'une pharmacie}
{Image UML à insérer : recherche mobile, appel API, filtrage par texte ou proximité, calcul de distance et retour des résultats.}
{fig:sequence-search-placeholder}

\placeholderfigure
{Diagramme de composants}
{Image UML à insérer : application mobile, dashboard administrateur, API FastAPI, services métier, ORM SQLAlchemy et base de données.}
{fig:component-diagram-placeholder}

\placeholderfigure
{Diagramme de déploiement}
{Image UML à insérer : client mobile, navigateur administrateur, serveur backend, base de données et services de configuration.}
{fig:deployment-diagram-placeholder}

\section{Sécurité et contrôle d'accès}

La sécurité a été traitée au niveau du backend, car les règles d'accès ne doivent pas dépendre uniquement
des interfaces clientes.

\subsection{Authentification JWT}

Après une connexion réussie, le backend génère un jeton d'accès et un jeton de rafraîchissement. Le
jeton d'accès sert à appeler les routes protégées. Le jeton de rafraîchissement permet de renouveler la
session lorsque le jeton d'accès expire.

\subsection{Contrôle d'accès basé sur les rôles}

Le système distingue plusieurs rôles administratifs, notamment administrateur, super administrateur et
assistant. Ces rôles limitent les actions accessibles. Les assistants peuvent aussi être liés à un périmètre
régional afin de limiter leurs interventions aux zones autorisées.

\subsection{Hachage des mots de passe}

Les mots de passe ne sont pas stockés en clair. Le backend les hache avant l'enregistrement et compare
les valeurs hachées pendant la connexion.

\subsection{Validation des API}

Les données envoyées vers l'API sont validées avec Pydantic et avec les services métier. Cette validation
concerne les champs obligatoires, les coordonnées GPS, les dates, les fichiers CSV, les limites de
pagination et les autorisations.

\subsection{Variables d'environnement}

Les paramètres sensibles sont placés dans des variables d'environnement. Cela concerne notamment la
chaîne de connexion à la base de données, la clé de signature JWT, la durée des jetons, les origines CORS,
les hôtes autorisés et la configuration du cache.

\section{Justification des choix techniques}

\begin{longtable}{L{4cm}L{10cm}}
    \caption{Justification des choix techniques}\\
    \toprule
    \textbf{Choix} & \textbf{Justification} \\
    \midrule
    \endfirsthead
    \toprule
    \textbf{Choix} & \textbf{Justification} \\
    \midrule
    \endhead
    FastAPI & Framework Python adapté aux API REST, avec validation Pydantic et documentation OpenAPI automatique. \\
    React et Vite & Ensemble adapté à la création d'une interface d'administration rapide, modulaire et maintenable. \\
    Expo et React Native & Permettent de développer une application mobile avec une base de code JavaScript et un accès simplifié aux fonctionnalités natives. \\
    SQLAlchemy & ORM permettant de structurer l'accès aux données tout en conservant un contrôle sur les modèles et les requêtes. \\
    SQLite et PostgreSQL & SQLite facilite le développement local, tandis que PostgreSQL constitue une option plus adaptée à un environnement configuré. \\
    JWT & Mécanisme adapté à la sécurisation des API et à la séparation entre jeton d'accès et jeton de rafraîchissement. \\
    CSV & Format simple pour importer des volumes de données de pharmacies, gardes et médicaments. \\
    \bottomrule
\end{longtable}

\section{Conclusion}

Ce chapitre a présenté la conception de PharmacieConnect. L'architecture retenue sépare le backend,
l'administration web, l'application mobile et la base de données. Cette organisation facilite le
développement, les tests et les évolutions futures. Le chapitre suivant présente la réalisation technique
et la validation de la solution.
